\paragraph{Results: one clear success, one entangled edit, two misses.}
\begin{itemize}
\item \textbf{\texttt{eli5} worked.} It cut the Flesch--Kincaid grade of answers from \EliNormalFK{} to \EliModeFK{} and halved their length. That is close to the readability of simply asking for a five-year-old explanation (\EliInstrFK), though the prompted answers were shorter still.
\item \textbf{\texttt{plain-math-add} was entangled with digits.} It removed LaTeX from more answers than the instruction did (\MathAddModeTex{} vs.\ \MathInstrTex{} of 20 still had LaTeX). But only \MathAddModeNum/20 answers still contained the correct number, against 20/20 for the instruction. The edit dropped or mangled digits: ``46 km/h multiplied by \textvisiblespace{} hours''. In the contrastive examples, numbers appear \emph{inside} the LaTeX, so the ``LaTeX direction'' also encoded part of how numbers are written.
\item \textbf{The weight edit \texttt{plain-math} did nothing on held-out prompts.} This is a rank-1 ablation of that direction from all 32 MLP layers, and it left LaTeX in \MathModeTex/20 answers. Removing a direction only from the MLP writes leaves the attention path and the residual stream free to carry it. Weight-Space Programming succeeded with a similar edit after a search over which columns to edit~\cite{takehi2026wsp}. Building one from 8 examples in seconds did not.
\item \textbf{\texttt{bullets} failed.} Normal Qwen already answered \BulNormalBul/20 of these how-to prompts mostly in bullets. The edit produced \BulModeBul/20, and the instruction \BulInstrBul/20.
\end{itemize}
For style, prompting was as good as or better than a quickly built edit. What the loop does show is that a frontier model can turn a request into a live, measurable edit in seconds, and that \texttt{try\_mode} catches broken strengths. Building edits that beat prompting will need better examples, held-out checks inside the loop, and probably edits on the attention path as well.
